\chapter{2020年江苏卷}

\section{填空题}

\begin{problem}
已知集合 \(A=\{-1,0,1,2\}\)，\(B=\{0,2,3\}\)，则 \(A\cap B=\)\fillinblank{}.
\end{problem}

\begin{answer}
\(\{0,2\}\)
\end{answer}

\begin{solution}
交集中的元素必须同时属于 \(A\) 和 \(B\)，所以 \(A\cap B=\{0,2\}\).
\end{solution}

\begin{problem}
已知 \(\i\) 是虚数单位，则复数 \(z=(1+\i)(2-\i)\) 的实部是\fillinblank{}.
\end{problem}

\begin{answer}
\(3\)
\end{answer}

\begin{solution}
\(z=(1+\i)(2-\i)=2-\i+2\i-\i^2=3+\i\)，所以复数 \(z\) 的实部为 \(3\).
\end{solution}

\begin{problem}
已知一组数据 \(4\)，\(2a\)，\(3-a\)，\(5\)，\(6\) 的平均数为 \(4\)，则 \(a\) 的值是\fillinblank{}.
\end{problem}

\begin{answer}
\(2\)
\end{answer}

\begin{solution}
由平均数为 \(4\)，得 \(4+2a+(3-a)+5+6=5\times4\)，即 \(a+18=20\)，所以 \(a=2\).
\end{solution}

\begin{problem}
将一颗质地均匀的正方体骰子先后抛掷 \(2\) 次，观察向上的点数，则点数和为 \(5\) 的概率是\fillinblank{}.
\end{problem}

\begin{answer}
\(\dfrac19\)
\end{answer}

\begin{solution}
两次抛掷的所有结果共有 \(6\times6=36\) 种，点数和为 \(5\) 的结果为 \((1,4)\)，\((2,3)\)，\((3,2)\)，\((4,1)\)，共 \(4\) 种，所以所求概率为 \(\dfrac4{36}=\dfrac19\).
\end{solution}

\begin{problem}
如图是一个算法流程图，若输出 \(y\) 的值为 \(-2\)，则输入 \(x\) 的值是

\bitmapfigure[max width=0.42\linewidth]{img/2020/jiangsu/q5_fig1.png}

\fillinblank{}.
\end{problem}

\begin{answer}
\(-3\)
\end{answer}

\begin{solution}
当 \(x>0\) 时，输出 \(y=2^x>0\)，不可能等于 \(-2\). 因而应取 \(x\leq0\) 的分支，此时 \(y=x+1\). 由 \(x+1=-2\)，得 \(x=-3\)，且满足 \(x\leq0\).
\end{solution}

\begin{problem}
在平面直角坐标系 \(xOy\) 中，若双曲线 \(\dfrac{x^2}{a^2}-\dfrac{y^2}{5}=1\)（\(a>0\)）的一条渐近线方程为 \(y=\dfrac{\sqrt5}{2}x\)，则该双曲线的离心率是\fillinblank{}.
\end{problem}

\begin{answer}
\(\dfrac32\)
\end{answer}

\begin{solution}
双曲线的渐近线方程为 \(y=\pm\dfrac{\sqrt5}{a}x\)，所以 \(\dfrac{\sqrt5}{a}=\dfrac{\sqrt5}{2}\)，得 \(a=2\). 由 \(c^2=a^2+5\)，得 \(c=3\)，故离心率 \(e=\dfrac ca=\dfrac32\).
\end{solution}

\begin{problem}
已知 \(y=f(x)\) 是奇函数，当 \(x\ge0\) 时，\(f(x)=x^{\frac23}\)，则 \(f(-8)\) 的值是\fillinblank{}.
\end{problem}

\begin{answer}
\(-4\)
\end{answer}

\begin{solution}
由奇函数的性质，\(f(-8)=-f(8)\). 又因为 \(f(8)=8^{\frac23}=\left(\sqrt[3]{8}\right)^2=4\)，所以 \(f(-8)=-4\).
\end{solution}

\begin{problem}
已知 \(\sin^2\!\left(\dfrac\pi4+\alpha\right)=\dfrac23\)，则 \(\sin2\alpha\) 的值是\fillinblank{}.
\end{problem}

\begin{answer}
\(\dfrac13\)
\end{answer}

\begin{solution}
利用 \(\sin^2u=\dfrac{1-\cos2u}{2}\)，取 \(u=\dfrac\pi4+\alpha\)，得
\[
\dfrac23=\dfrac{1-\cos\left(\dfrac\pi2+2\alpha\right)}2=\dfrac{1+\sin2\alpha}2.
\]
因此 \(\sin2\alpha=\dfrac13\).
\end{solution}

\begin{problem}
如图，六角螺帽毛坯是由一个正六棱柱挖去一个圆柱所构成的. 已知螺帽的底面正六边形边长为 \(2\mathrm{~cm}\)，高为 \(2\mathrm{~cm}\)，内孔半径为 \(0.5\mathrm{~cm}\)，则此六角螺帽毛坯的体积是

\bitmapfigure[max width=0.42\linewidth]{img/2020/jiangsu/q9_fig1.png}

\fillinblank{}\(\mathrm{cm}^3\).
\end{problem}

\begin{answer}
\(12\sqrt3-\dfrac\pi2\)
\end{answer}

\begin{solution}
正六边形可分成 \(6\) 个边长为 \(2\) 的正三角形，所以底面积为 \(6\times\dfrac{\sqrt3}{4}\times2^2=6\sqrt3\). 正六棱柱体积为 \(6\sqrt3\times2=12\sqrt3\). 内孔圆柱体积为 \(\pi\times0.5^2\times2=\dfrac\pi2\). 因而毛坯体积为 \(12\sqrt3-\dfrac\pi2\ \mathrm{cm}^3\).
\end{solution}

\begin{problem}
将函数 \(y=3\sin\!\left(2x+\dfrac\pi4\right)\) 的图象向右平移 \(\dfrac\pi6\) 个单位长度，则平移后的图象中与 \(y\) 轴最近的对称轴的方程是\fillinblank{}.
\end{problem}

\begin{answer}
\(x=-\dfrac{5\pi}{24}\)
\end{answer}

\begin{solution}
向右平移后函数变为 \(y=3\sin\!\left(2\left(x-\dfrac\pi6\right)+\dfrac\pi4\right)=3\sin\!\left(2x-\dfrac\pi{12}\right)\). 正弦函数的对称轴满足
\[
2x-\dfrac\pi{12}=\dfrac\pi2+k\pi,\qquad k\in\Z,
\]
即 \(x=\dfrac{7\pi}{24}+\dfrac{k\pi}{2}\). 取 \(k=-1\) 时得到距离 \(y\) 轴最近的对称轴 \(x=-\dfrac{5\pi}{24}\).
\end{solution}

\begin{problem}
设 \(\{a_n\}\) 是公差为 \(d\) 的等差数列，\(\{b_n\}\) 是公比为 \(q\) 的等比数列. 已知数列 \(\{a_n+b_n\}\) 的前 \(n\) 项和 \(S_n=n^2-n+2^n-1\)（\(n\in\N^*\)），则 \(d+q\) 的值是\fillinblank{}.
\end{problem}

\begin{answer}
\(4\)
\end{answer}

\begin{solution}
令 \(c_n=a_n+b_n\). 由 \(S_n-S_{n-1}=c_n\)（并直接验证 \(n=1\)），得
\[
c_n=2n-2+2^{n-1}.
\]
于是 \(c_1=1\)，且 \(c_2-c_1=3\)，\(c_3-c_2=4\)，\(c_4-c_3=6\). 设 \(b_1=B\)，则
\[
c_{n+1}-c_n=d+B(q-1)q^{n-1}.
\]
相邻两项作差，得到
\[
1=B(q-1)^2,\qquad 2=Bq(q-1)^2.
\]
所以 \(q=2\). 再由 \(d+B(q-1)=3\) 和 \(d+2B(q-1)=4\)，得 \(B=1\)，\(d=2\). 因而 \(d+q=4\).
\end{solution}

\begin{problem}
已知 \(5x^2y^2+y^4=1\)（\(x,y\in\R\)），则 \(x^2+y^2\) 的最小值是\fillinblank{}.
\end{problem}

\begin{answer}
\(\dfrac45\)
\end{answer}

\begin{solution}
令 \(u=y^2\). 由题意 \(u>0\)，且 \(u\leq1\)，并有 \(x^2=\dfrac{1-u^2}{5u}\). 因而
\[
x^2+y^2=\dfrac{1-u^2}{5u}+u=\dfrac1{5u}+\dfrac45u\geq\dfrac45,
\]
其中等号当且仅当 \(\dfrac1{5u}=\dfrac45u\)，即 \(u=\dfrac12\). 所以最小值为 \(\dfrac45\).
\end{solution}

\begin{problem}
在 \(\triangle ABC\) 中，\(AB=4\)，\(AC=3\)，\(\angle BAC=90^\circ\). \(D\) 在边 \(BC\) 上，延长 \(AD\) 到 \(P\)，使得 \(AP=9\). 若 \(\overrightarrow{PA}=m\overrightarrow{PB}+\left(\dfrac32-m\right)\overrightarrow{PC}\)（\(m\) 为常数），则 \(CD\) 的长度是

\bitmapfigure[max width=0.42\linewidth]{img/2020/jiangsu/q13_fig1.png}

\fillinblank{}.
\end{problem}

\begin{answer}
\(\dfrac{18}{5}\)
\end{answer}

\begin{solution}
因为 \(P,D,A\) 共线且 \(P\) 在 \(D\) 的外侧，可设 \(\overrightarrow{PA}=\lambda\overrightarrow{PD}\)，其中 \(\lambda>0\). 代入已知向量关系，得
\[
\overrightarrow{PD}=\dfrac{m}{\lambda}\overrightarrow{PB}+\dfrac{\frac32-m}{\lambda}\overrightarrow{PC}.
\]
右侧是点 \(B,C\) 的仿射组合，而其结果为点 \(D\) 的位置向量；由 \(D\) 在直线 \(BC\) 上，两个系数之和应为 \(1\)，所以 \(\dfrac{3}{2\lambda}=1\)，得 \(\lambda=\dfrac32\). 因而 \(PD=\dfrac{AP}{\lambda}=6\)，从而 \(AD=AP-PD=3\). 又 \(BC=\sqrt{AB^2+AC^2}=5\).

设 \(CD=x\)，则 \(BD=5-x\)，令 \(\angle CDA=\theta\)，于是 \(\angle BDA=\pi-\theta\). 在 \(\triangle ACD\) 和 \(\triangle ABD\) 中分别用余弦定理，得
\[
\cos\theta=\dfrac{AD^2+CD^2-AC^2}{2AD\cdot CD}=\dfrac{x}{6},
\]
\[
\cos(\pi-\theta)=\dfrac{AD^2+BD^2-AB^2}{2AD\cdot BD}=\dfrac{(5-x)^2-7}{6(5-x)}.
\]
由于 \(\cos(\pi-\theta)=-\cos\theta\)，所以
\[
\dfrac{x}{6}+\dfrac{(5-x)^2-7}{6(5-x)}=0.
\]
化简得 \(18-5x=0\)，故 \(CD=x=\dfrac{18}{5}\).
\end{solution}

\begin{problem}
在平面直角坐标系 \(xOy\) 中，已知 \(P\!\left(\dfrac{\sqrt3}{2},0\right)\)，\(A\)，\(B\) 是圆 \(C:x^2+\left(y-\dfrac12\right)^2=36\) 上的两个动点，满足 \(PA=PB\)，则 \(\triangle PAB\) 面积的最大值是\fillinblank{}.
\end{problem}

\begin{answer}
\(10\sqrt5\)
\end{answer}

\begin{solution}
记圆心为 \(O_1\left(0,\dfrac12\right)\)，则 \(O_1P=1\). 由 \(PA=PB\)，点 \(P\) 在弦 \(AB\) 的垂直平分线上，所以 \(O_1P\perp AB\). 设圆心到弦 \(AB\) 的距离为 \(d\)，其中 \(0\leq d<6\)，则 \(AB=2\sqrt{36-d^2}\). 对于固定的 \(d\)，取弦在圆心的另一侧时，点 \(P\) 到弦的距离为 \(d+1\)，这给出最大可能的三角形面积. 因此
\[
S_{\triangle PAB}\leq\sqrt{36-d^2}(d+1).
\]
令 \(Y=(36-d^2)(d+1)^2\)，则
\[
Y'=2(d+1)(36-2d^2-d).
\]
在 \([0,6)\) 内，\(Y'=0\) 的唯一解为 \(d=4\)，且函数先增后减，所以面积最大值为
\[
\sqrt{36-4^2}(4+1)=\sqrt{20}\times5=10\sqrt5.
\]
\end{solution}

\section{解答题}

\begin{problem}
在三棱柱 \(ABC-A_1B_1C_1\) 中，\(AB\perp AC\)，\(B_1C\perp\) 平面 \(ABC\)，\(E\)，\(F\) 分别是 \(AC\)，\(B_1C\) 的中点.

\begin{enumerate}
\item 求证：\(EF\parallel\) 平面 \(AB_1C_1\)；
\item 求证：平面 \(AB_1C\perp\) 平面 \(ABB_1\).
\end{enumerate}

\bitmapfigure[max width=0.42\linewidth]{img/2020/jiangsu/q15_fig1.png}
\end{problem}

\begin{solution}
\begin{enumerate}
\item 在三角形 \(ACB_1\) 中，\(E\)，\(F\) 分别是 \(AC\)，\(CB_1\) 的中点，所以由中位线定理得 \(EF\parallel AB_1\). 又因为 \(AB_1\subset\) 平面 \(AB_1C_1\)，故 \(EF\parallel\) 平面 \(AB_1C_1\).
\item 由 \(B_1C\perp\) 平面 \(ABC\)，且 \(AB\subset\) 平面 \(ABC\)，得 \(B_1C\perp AB\). 已知 \(AB\perp AC\)，而 \(AC\) 与 \(B_1C\) 是平面 \(AB_1C\) 内相交的两条直线，所以 \(AB\perp\) 平面 \(AB_1C\). 又 \(AB\subset\) 平面 \(ABB_1\)，因此平面 \(AB_1C\perp\) 平面 \(ABB_1\).
\end{enumerate}
\end{solution}

\begin{problem}
在 \(\triangle ABC\) 中，角 \(A\)，\(B\)，\(C\) 的对边分别为 \(a\)，\(b\)，\(c\)，已知 \(a=3\)，\(c=\sqrt2\)，\(B=45^\circ\).

\begin{enumerate}
\item 求 \(\sin C\) 的值；
\item 在边 \(BC\) 上取一点 \(D\)，使得 \(\cos\angle ADC=-\dfrac45\)，求 \(\tan\angle DAC\) 的值.
\end{enumerate}

\bitmapfigure[max width=0.42\linewidth]{img/2020/jiangsu/q16_fig1.png}
\end{problem}

\begin{solution}
\begin{enumerate}
\item 由余弦定理，\(b^2=a^2+c^2-2ac\cos B=9+2-2\times3\times\sqrt2\times\dfrac{\sqrt2}{2}=5\)，所以 \(b=\sqrt5\). 再由正弦定理，得
\[
\sin C=\dfrac{c\sin B}{b}=\dfrac{\sqrt2\times\frac{\sqrt2}{2}}{\sqrt5}=\dfrac{\sqrt5}{5}.
\]
\item 因为 \(\cos\angle ADC=-\dfrac45\)，且 \(\angle ADC\in(0,\pi)\)，所以 \(\sin\angle ADC=\dfrac35\). 由 \(\sin C=\dfrac{\sqrt5}{5}\) 且 \(B=45^\circ\)，可知 \(C<90^\circ\)，从而 \(\cos C=\dfrac{2\sqrt5}{5}\). 在三角形 \(ADC\) 中，\(\angle DAC=\pi-\angle ADC-C\)，故
\[
\sin\angle DAC=\sin(\angle ADC+C)=\dfrac35\cdot\dfrac{2\sqrt5}{5}-\dfrac45\cdot\dfrac{\sqrt5}{5}=\dfrac{2\sqrt5}{25},
\]
\[
\cos\angle DAC=\dfrac{11\sqrt5}{25}.
\]
因此 \(\tan\angle DAC=\dfrac{\sin\angle DAC}{\cos\angle DAC}=\dfrac2{11}\).
\end{enumerate}
\end{solution}

\begin{problem}
某地准备在山谷中建一座桥梁，桥址位置的竖直截面图如图所示：谷底 \(O\) 在水平线 \(MN\) 上，桥 \(AB\) 与 \(MN\) 平行，\(OO'\) 为铅垂线（\(O'\) 在 \(AB\) 上）. 经测量，左侧曲线 \(AO\) 上任一点 \(D\) 到 \(MN\) 的距离 \(h_1\)（\(\mathrm{~m}\)）与 \(D\) 到 \(OO'\) 的距离 \(a\)（\(\mathrm{~m}\)）之间满足关系式 \(h_1=\dfrac{1}{40}a^2\)；右侧曲线 \(BO\) 上任一点 \(F\) 到 \(MN\) 的距离 \(h_2\)（\(\mathrm{~m}\)）与 \(F\) 到 \(OO'\) 的距离 \(b\)（\(\mathrm{~m}\)）之间满足关系式 \(h_2=-\dfrac{1}{800}b^3+6b\). 已知点 \(B\) 到 \(OO'\) 的距离为 \(40\mathrm{~m}\).

\begin{enumerate}
\item 求桥 \(AB\) 的长度；
\item 计划在谷底两侧建造平行于 \(OO'\) 的桥墩 \(CD\) 和 \(EF\)，且 \(CE\) 为 \(80\mathrm{~m}\)，其中 \(C\)，\(E\) 在 \(AB\) 上（不包括端点）. 桥墩 \(EF\) 每米造价 \(k\)（万元），桥墩 \(CD\) 每米造价 \(\dfrac32k\)（万元）（\(k>0\)）. 问 \(O'E\) 为多少米时，桥墩 \(CD\) 与 \(EF\) 的总造价最低？
\end{enumerate}

\bitmapfigure[max width=0.42\linewidth]{img/2020/jiangsu/q17_fig1.png}
\end{problem}

\begin{solution}
\begin{enumerate}
\item 设 \(O'A=a_A\). 因为 \(B\) 到 \(OO'\) 的距离为 \(40\mathrm{~m}\)，所以桥面高度为
\[
160=-\dfrac{1}{800}\times40^3+6\times40.
\]
由左侧曲线关系式，\(\dfrac1{40}a_A^2=160\)，得 \(a_A=80\). 因而 \(AB=80+40=120\ \mathrm{m}\).
\item 设 \(O'E=x\). 由图形位置及 \(C,E\) 不为端点，得 \(0<x<40\)，且 \(O'C=80-x\). 桥墩长度分别为
\[
EF=160+\dfrac{x^3}{800}-6x,
\qquad
CD=160-\dfrac{(80-x)^2}{40}.
\]
总造价除以正数 \(k\) 后为
\[
T(x)=160+\dfrac{x^3}{800}-6x+\dfrac32\left(160-\dfrac{(80-x)^2}{40}\right).
\]
所以
\[
T'(x)=\dfrac{3x^2}{800}-6+\dfrac{3(80-x)}{40}=\dfrac{3x(x-20)}{800}.
\]
当 \(0<x<20\) 时 \(T'(x)<0\)，当 \(20<x<40\) 时 \(T'(x)>0\)，故总造价在 \(x=20\) 时取得最小值，即 \(O'E=20\ \mathrm{m}\).
\end{enumerate}
\end{solution}

\begin{problem}
在平面直角坐标系 \(xOy\) 中，已知椭圆 \(E:\dfrac{x^2}{4}+\dfrac{y^2}{3}=1\) 的左、右焦点分别为 \(F_1\)，\(F_2\)，点 \(A\) 在椭圆 \(E\) 上且在第一象限内，\(AF_2\perp F_1F_2\)，直线 \(AF_1\) 与椭圆 \(E\) 相交于另一点 \(B\).

\begin{enumerate}
\item 求 \(\triangle AF_1F_2\) 的周长；
\item 在 \(x\) 轴上任取一点 \(P\)，直线 \(AP\) 与椭圆 \(E\) 的右准线相交于点 \(Q\)，求 \(\overrightarrow{OP}\cdot\overrightarrow{QP}\) 的最小值；
\item 设点 \(M\) 在椭圆 \(E\) 上，记 \(\triangle OAB\) 与 \(\triangle MAB\) 的面积分别为 \(S_1\)，\(S_2\)，若 \(S_2=3S_1\)，求点 \(M\) 的坐标.
\end{enumerate}

\bitmapfigure[max width=0.42\linewidth]{img/2020/jiangsu/q18_fig1.png}
\end{problem}

\begin{solution}
\begin{enumerate}
\item 椭圆的长半轴为 \(2\)，焦距的一半为 \(c=\sqrt{4-3}=1\)，所以 \(F_1=(-1,0)\)，\(F_2=(1,0)\). 由 \(AF_2\perp F_1F_2\)，得 \(x_A=1\). 代入椭圆方程并取第一象限的点，得 \(A=\left(1,\dfrac32\right)\). 由椭圆定义，\(AF_1+AF_2=4\)，故三角形周长为 \(4+F_1F_2=4+2=6\).
\item 设 \(P=(p,0)\). 椭圆的右准线为 \(x=4\). 直线 \(AP\) 与准线交于 \(Q=(4,y_Q)\)，因此
\[
\overrightarrow{OP}\cdot\overrightarrow{QP}=(p,0)\cdot(p-4,-y_Q)=p(p-4)=(p-2)^2-4\geq-4.
\]
取 \(p=2\) 时直线确与准线相交，所以最小值为 \(-4\).
\item 直线 \(AB\) 即直线 \(AF_1\)，方程为 \(3x-4y+3=0\). 原点到此直线的距离为 \(\dfrac35\). 由两个三角形有公共底边 \(AB\)，且 \(S_2=3S_1\)，得点 \(M\) 到直线 \(AB\) 的距离为 \(\dfrac95\). 因而
\[
\left|3x_M-4y_M+3\right|=9.
\]
当 \(3x_M-4y_M+3=9\) 时，联立椭圆方程得
\[
7x_M^2-12x_M-4=0,
\]
所以 \(M=(2,0)\) 或 \(M=\left(-\dfrac27,-\dfrac{12}{7}\right)\). 当 \(3x_M-4y_M+3=-9\) 时，联立得到 \(7x_M^2+24x_M+32=0\)，其判别式小于零，无实数解. 故点 \(M\) 的坐标为 \((2,0)\) 或 \(\left(-\dfrac27,-\dfrac{12}{7}\right)\).
\end{enumerate}
\end{solution}

\begin{problem}
已知关于 \(x\) 的函数 \(y=f(x)\)，\(y=g(x)\) 与 \(h(x)=kx+b\)（\(k,b\in\R\)）在区间 \(D\) 上恒有 \(f(x)\ge h(x)\ge g(x)\).

\begin{enumerate}
\item 若 \(f(x)=x^2+2x\)，\(g(x)=-x^2+2x\)，\(D=(-\infty,+\infty)\)，求 \(h(x)\) 的表达式；
\item 若 \(f(x)=x^2-x+1\)，\(g(x)=k\ln x\)，\(h(x)=kx-k\)，\(D=(0,+\infty)\)，求 \(k\) 的取值范围；
\item 若 \(f(x)=x^4-2x^2\)，\(g(x)=4x^2-8\)，\(h(x)=4(t^3-t)x-3t^4+2t^2\)，\(0<|t|\le\sqrt2\)，\(D=[m,n]\subseteq[-\sqrt2,\sqrt2]\)，证明：\(n-m\le\sqrt7\).
\end{enumerate}
\end{problem}

\begin{solution}
\begin{enumerate}
\item 因为 \(f(0)=g(0)=0\)，由恒有关系得 \(h(0)=0\)，所以 \(b=0\). 当 \(x>0\) 时，\(-x^2+2x\leq kx\leq x^2+2x\)，即 \(2-x\leq k\leq2+x\). 令 \(x\to0^+\)，得 \(k=2\). 直接代入可得 \(x^2+2x\geq2x\geq-x^2+2x\)，所以 \(h(x)=2x\).
\item 由 \(x-1\geq\ln x\)（\(x>0\)）可知
\[
h(x)-g(x)=k(x-1-\ln x).
\]
当 \(x\ne1\) 时括号内严格为正，所以必须且只需 \(k\geq0\). 另一方面
\[
f(x)-h(x)=x^2-(k+1)x+(k+1).
\]
在 \(k\geq0\) 时，这是开口向上的二次函数；要使其在 \((0,+\infty)\) 恒非负，需 \((k+1)^2-4(k+1)=(k+1)(k-3)\leq0\)，即 \(0\leq k\leq3\). 这个范围内两项不等式均成立，故 \(k\in[0,3]\).
\item 先因式分解：
\[
f(x)-h(x)=(x-t)^2\bigl((x+t)^2+2t^2-2\bigr).
\]
又
\[
g(x)-h(x)=4x^2-4(t^3-t)x+3t^4-2t^2-8.
\]
令 \(u=t^2\)，并设上式的两个根为 \(x_1\leq x_2\). 因为 \(D\) 非空且在 \(D\) 上有 \(h\geq g\)，所以根为实数，且 \(D\subseteq[x_1,x_2]\). 由韦达定理，
\[
x_1+x_2=t^3-t,
\qquad
(x_2-x_1)^2=t^6-5t^4+3t^2+8=P(u),
\]
其中 \(P(u)=u^3-5u^2+3u+8\).

若 \(1\leq u\leq2\)，则 \(P'(u)=(3u-1)(u-3)<0\)，从而 \(P(u)\leq P(1)=7\). 因此 \(n-m\leq x_2-x_1\leq\sqrt7\).

下面设 \(0<u<1\). 由于 \(f,g\) 都是偶函数，且 \(h_{-t}(x)=h_t(-x)\)，把 \(t<0\) 与区间 \([m,n]\) 同时关于原点对称不改变区间长度，所以只需考虑 \(t>0\). 令
\[
c=\dfrac{x_1+x_2}{2}=\dfrac{t(u-1)}2,\qquad w=\dfrac{x_2-x_1}{2}=\dfrac{\sqrt{P(u)}}2,qquad r=\sqrt{2-2u}.
\]
由 \(f(x)-h(x)\geq0\)，除去可能的孤立点 \(x=t\) 外，区间只能位于 \(x\leq-t-r\) 或 \(x\geq-t+r\). 此外
\[
g(-t)-h(-t)=7u^2-2u-8<0,
\]
所以 \(-t\) 位于 \((x_1,x_2)\) 内. 因而非退化的区间 \(D\) 只能落在上述两个允许部分之一，其中右侧允许部分的长度至多为
\[
x_2-(-t+r)=w+\dfrac{t(1+u)}2-r,
\]
左侧允许部分的长度至多为 \(w-\dfrac{t(1+u)}2-r\)，故
\[
n-m\leq w+\dfrac{t(1+u)}2-r.
\]
又 \(P'(u)\) 在 \((0,1)\) 内的唯一零点为 \(u=\dfrac13\)，且 \(P\left(\dfrac13\right)=\dfrac{229}{27}<9\)，所以 \(w<\dfrac32\). 由于 \(0<t<1\)、\(0<\dfrac{1+u}{2}<1\) 且 \(r>0\)，因此
\[
n-m<\dfrac32+1=\dfrac52<\sqrt7.
\]
综上，始终有 \(n-m\leq\sqrt7\).
\end{enumerate}
\end{solution}

\begin{problem}
已知数列 \(\{a_n\}\)（\(n\in\N^*\)）的首项 \(a_1=1\)，前 \(n\) 项和为 \(S_n\). 设 \(\lambda\) 与 \(k\) 是常数，若对一切正整数 \(n\)，均有 \(S_{n+1}^{1/k}-S_n^{1/k}=\lambda a_{n+1}^{1/k}\) 成立，则称此数列为“\(\lambda\sim k\)”数列.

\begin{enumerate}
\item 若等差数列 \(\{a_n\}\) 是“\(\lambda\sim1\)”数列，求 \(\lambda\) 的值；
\item 若数列 \(\{a_n\}\) 是“\(\dfrac{\sqrt3}{3}\sim2\)”数列，且 \(a_n>0\)，求数列 \(\{a_n\}\) 的通项公式；
\item 对于给定的 \(\lambda\)，是否存在三个不同的数列 \(\{a_n\}\) 为“\(\lambda\sim3\)”数列，且 \(a_n\ge0\)？若存在，求 \(\lambda\) 的取值范围；若不存在，说明理由.
\end{enumerate}
\end{problem}

\begin{solution}
\begin{enumerate}
\item 当 \(k=1\) 时，定义式化为 \(a_{n+1}=\lambda a_{n+1}\). 等差数列不可能从第二项起所有项都为零，因为 \(a_1=1\)，所以存在某个 \(n\) 使 \(a_{n+1}\ne0\)，从而 \(\lambda=1\).
\item 由 \(a_n>0\)，有 \(S_{n+1}>S_n>0\). 令 \(U=\sqrt{S_{n+1}}\)，\(V=\sqrt{S_n}\)，则
\[
U-V=\dfrac{\sqrt3}{3}\sqrt{U^2-V^2}.
\]
因 \(U>V\)，平方并约去正因子 \(U-V\)，得 \(U-V=\dfrac{U+V}{3}\)，所以 \(U=2V\)，即 \(S_{n+1}=4S_n\). 由 \(S_1=1\)，得 \(S_n=4^{n-1}\). 因而
\[
a_1=1,\qquad a_n=S_n-S_{n-1}=3\cdot4^{n-2}\quad(n\ge2).
\]
\item 若某一步 \(a_{n+1}>0\)，令 \(r=\dfrac{S_{n+1}^{1/3}}{S_n^{1/3}}>1\). 由定义式并立方化简，得
\[
(r-1)^2=\lambda^3(r^2+r+1).
\]
左边与右侧的比值为 \(\dfrac{(r-1)^2}{r^2+r+1}<1\)，所以发生正增量必须有 \(0<\lambda<1\). 当 \(0<\lambda<1\) 时，令 \(L=\lambda^3\)，方程化为
\[
(1-L)r^2-(2+L)r+(1-L)=0.
\]
其判别式为 \(3L(4-L)>0\)，两根乘积为 \(1\)，且两根和为正，因此存在唯一的根 \(r>1\). 记 \(R=r^3>1\). 每一步都可以选择保持 \(S_{n+1}=S_n\)，也可以选择跳到 \(S_{n+1}=RS_n\)，均满足定义式.

例如，以下三个数列均满足条件：第一列只有 \(a_1=1\) 非零；第二列取 \(a_2=R-1\) 后其余为零；第三列取 \(a_2=R-1\)，\(a_3=R^2-R\) 后其余为零. 它们彼此不同且各项非负. 因此存在三个不同数列的充要条件为 \(0<\lambda<1\).
\end{enumerate}
\end{solution}

\section{附加题：解答题}

\begin{problem}[A]
平面上点 \(A(2,-1)\) 在矩阵 \(M=\begin{bmatrix}a&1\\-1&b\end{bmatrix}\) 对应的变换作用下得到点 \(B(3,-4)\).

\begin{enumerate}
\item 求实数 \(a\)，\(b\) 的值；
\item 求矩阵 \(M\) 的逆矩阵 \(M^{-1}\).
\end{enumerate}
\end{problem}

\begin{solution}
\begin{enumerate}
\item 由矩阵变换关系，
\[
\begin{bmatrix}a&1\\-1&b\end{bmatrix}\begin{bmatrix}2\\-1\end{bmatrix}=\begin{bmatrix}3\\-4\end{bmatrix},
\]
所以 \(2a-1=3\)，\(-2-b=-4\)，解得 \(a=2\)，\(b=2\).
\item 此时 \(M=\begin{bmatrix}2&1\\-1&2\end{bmatrix}\)，其行列式为 \(5\)，故
\[
M^{-1}=\dfrac15\begin{bmatrix}2&-1\\1&2\end{bmatrix}.
\]
\end{enumerate}
\end{solution}

\reuseproblemnumber
\begin{problem}[B]
在极坐标系中，已知点 \(A\!\left(\rho_1,\dfrac\pi3\right)\) 在直线 \(l:\rho\cos\theta=2\) 上，点 \(B\!\left(\rho_2,\dfrac\pi6\right)\) 在圆 \(C:\rho=4\sin\theta\) 上（其中 \(\rho\ge0\)，\(0\le\theta<2\pi\)）.

\begin{enumerate}
\item 求 \(\rho_1\)，\(\rho_2\) 的值；
\item 求出直线 \(l\) 与圆 \(C\) 的公共点的极坐标.
\end{enumerate}
\end{problem}

\begin{solution}
\begin{enumerate}
\item 将点 \(A\) 代入直线方程，得 \(\rho_1\cos\dfrac\pi3=2\)，所以 \(\rho_1=4\). 将点 \(B\) 代入圆方程，得 \(\rho_2=4\sin\dfrac\pi6=2\).
\item 公共点满足
\[
\rho\cos\theta=2,\qquad \rho=4\sin\theta.
\]
消去 \(\rho\)，得 \(4\sin\theta\cos\theta=2\)，即 \(\sin2\theta=1\). 在 \([0,2\pi)\) 内，\(\theta=\dfrac\pi4\) 或 \(\dfrac{5\pi}4\). 当 \(\theta=\dfrac{5\pi}4\) 时由直线方程得到 \(\rho=-2\sqrt2<0\)，不合题设；当 \(\theta=\dfrac\pi4\) 时 \(\rho=2\sqrt2\). 所以公共点的极坐标为 \(\left(2\sqrt2,\dfrac\pi4\right)\).
\end{enumerate}
\end{solution}

\reuseproblemnumber
\begin{problem}[C]
设 \(x\in\R\)，解不等式 \(2\lvert x+1\rvert+\lvert x\rvert\le4\).
\end{problem}

\begin{solution}
当 \(x<-1\) 时，不等式化为 \(-3x-2\le4\)，得 \(-2\le x<-1\)；当 \(-1\le x\le0\) 时，不等式化为 \(x+2\le4\)，该区间内所有 \(x\) 均满足；当 \(x>0\) 时，不等式化为 \(3x+2\le4\)，得 \(0<x\le\dfrac23\). 因而解集为 \(\left[-2,\dfrac23\right]\).
\end{solution}

\begin{problem}
在三棱锥 \(A-BCD\) 中，已知 \(CB=CD=\sqrt5\)，\(BD=2\)，\(O\) 为 \(BD\) 的中点，\(AO\perp\) 平面 \(BCD\)，\(AO=2\)，\(E\) 为 \(AC\) 的中点.

\begin{enumerate}
\item 求直线 \(AB\) 与 \(DE\) 所成角的余弦值；
\item 若点 \(F\) 在 \(BC\) 上，满足 \(BF=\dfrac14BC\)，设二面角 \(F-DE-C\) 的大小为 \(\theta\)，求 \(\sin\theta\) 的值.
\end{enumerate}

\bitmapfigure[max width=0.42\linewidth]{img/2020/jiangsu/q24_fig1.png}
\end{problem}

\begin{solution}
\begin{enumerate}
\item 取 \(O\) 为原点，分别以 \(OB\)，\(OC\)，\(OA\) 所在方向为 \(x\)，\(y\)，\(z\) 轴. 由 \(CB=CD\) 且 \(O\) 为 \(BD\) 的中点，得 \(CO\perp BD\)，可设
\[
A=(0,0,2),\quad B=(1,0,0),\quad C=(0,2,0),\quad D=(-1,0,0).
\]
于是 \(E=(0,1,1)\). 取 \(\overrightarrow{AB}=(1,0,-2)\)，\(\overrightarrow{DE}=(1,1,1)\)，则
\[
\cos\angle(AB,DE)=\dfrac{\left|\overrightarrow{AB}\cdot\overrightarrow{DE}\right|}{|\overrightarrow{AB}|\,|\overrightarrow{DE}|}=\dfrac{|-1|}{\sqrt5\sqrt3}=\dfrac{\sqrt{15}}{15}.
\]
\item 由 \(BF=\dfrac14BC\)，得 \(F=\left(\dfrac34,\dfrac12,0\right)\). 平面 \(CDE\) 内的两条方向向量可取 \(\overrightarrow{DC}=(1,2,0)\)，\(\overrightarrow{DE}=(1,1,1)\)，所以其一个法向量为 \(\bs n_1=(-2,1,1)\). 平面 \(FDE\) 内可取 \(\overrightarrow{DF}=\left(\dfrac74,\dfrac12,0\right)\)，\(\overrightarrow{DE}=(1,1,1)\)，所以其一个法向量为 \(\bs n_2=(2,-7,5)\). 因而
\[
\bs n_1\cdot\bs n_2=-6,\qquad |\bs n_1|=\sqrt6,\qquad |\bs n_2|=\sqrt{78}.
\]
二面角取两平面锐角，故
\[
\sin\theta=\sqrt{1-\left(\dfrac{-6}{\sqrt6\sqrt{78}}\right)^2}=\sqrt{1-\dfrac1{13}}=\dfrac{2\sqrt{39}}{13}.
\]
\end{enumerate}
\end{solution}

\begin{problem}
甲口袋中装有 \(2\) 个黑球和 \(1\) 个白球，乙口袋中装有 \(3\) 个白球. 现从甲，乙两口袋中各任取一个球交换放入另一口袋，重复 \(n\) 次这样的操作，记甲口袋中黑球个数为 \(X_n\)，恰有 \(2\) 个黑球的概率为 \(p_n\)，恰有 \(1\) 个黑球的概率为 \(q_n\).

\begin{enumerate}
\item 求 \(p_1\cdot q_1\) 和 \(p_2\cdot q_2\)；
\item 求 \(2p_n+q_n\) 与 \(2p_{n-1}+q_{n-1}\) 的递推关系式和 \(X_n\) 的数学期望 \(E(X_n)\)（用 \(n\) 表示）.
\end{enumerate}
\end{problem}

\begin{solution}
\begin{enumerate}
\item 第一次交换后，甲袋仍有 \(3\) 个球. 甲袋取黑球且乙袋取白球时，黑球数变为 \(1\)，所以 \(q_1=\dfrac23\)；甲袋取白球且乙袋取白球时，黑球数仍为 \(2\)，所以 \(p_1=\dfrac13\). 因而 \(p_1q_1=\dfrac29\).

第二次交换时，若交换前甲袋有 \(2\) 个黑球，则转移到 \(2\) 个黑球的概率为 \(\dfrac13\)，转移到 \(1\) 个黑球的概率为 \(\dfrac23\). 若交换前甲袋有 \(1\) 个黑球，则转移到 \(2\) 个黑球的概率为 \(\dfrac23\times\dfrac13=\dfrac29\)，保持 \(1\) 个黑球的概率为 \(\dfrac13\times\dfrac13+\dfrac23\times\dfrac23=\dfrac59\)，转移到 \(0\) 个黑球的概率为 \(\dfrac29\). 因此
\[
p_2=\dfrac13p_1+\dfrac29q_1=\dfrac7{27},\qquad q_2=\dfrac23p_1+\dfrac59q_1=\dfrac{16}{27},
\]
从而 \(p_2q_2=\dfrac{112}{729}\).
\item 还需考虑甲袋有 \(0\) 个黑球的状态. 对任意 \(n\ge1\)，由三种状态的转移概率得
\[
p_n=\dfrac13p_{n-1}+\dfrac29q_{n-1},
\]
\[
q_n=\dfrac23p_{n-1}+\dfrac59q_{n-1}+\dfrac23\left(1-p_{n-1}-q_{n-1}\right)=\dfrac23-\dfrac19q_{n-1}.
\]
所以
\[
2p_n+q_n=\dfrac13\left(2p_{n-1}+q_{n-1}\right)+\dfrac23.
\]
令 \(r_n=2p_n+q_n\)，则 \(r_1=\dfrac43\)，且 \(r_n-1=\dfrac13(r_{n-1}-1)\). 因而
\[
r_n=1+\dfrac1{3^n}.
\]
由于 \(X_n\) 取值为 \(0\)，\(1\)，\(2\)，故
\[
E(X_n)=2p_n+q_n=1+\dfrac1{3^n}.
\]
\end{enumerate}
\end{solution}
